%% CASS — KDD 2027 Research Track submission (anonymous)
\documentclass[sigconf,review,anonymous]{acmart}

\AtBeginDocument{%
  \providecommand\BibTeX{{Bib\TeX}}}

\setcopyright{none}
\copyrightyear{2027}
\acmYear{2027}
\acmConference[KDD '27]{The 33rd ACM SIGKDD Conference on Knowledge Discovery and Data Mining}{August 2027}{TBD}
\acmISBN{978-x-xxxx-xxxx-x/27/08}
\acmDOI{XXXXXXX.XXXXXXX}

\usepackage{booktabs}
\usepackage{multirow}
\usepackage{array}
\usepackage{graphicx}
\usepackage{amsmath,amssymb}
\usepackage{algorithm}
\usepackage{algpseudocode}

\newtheorem{proposition}{Proposition}
\newtheorem{assumption}{Assumption}

\newcommand{\method}{\textsc{CASS}}
\newcommand{\eg}{e.g.,}
\newcommand{\ie}{i.e.,}
\newcommand{\R}{\mathbb{R}}

\begin{document}

\title{\method: Training-Free Adaptation to Unseen Tasks by Mining and
Composing Skill Subspaces from LLM Activations}

\author{Anonymous Author(s)}
\affiliation{%
  \institution{Anonymous Institution}
  \city{}
  \country{}
}

\begin{abstract}
Large language models encode task information as manipulable structure in
their activation space: task vectors extracted from in-context demonstrations
causally trigger task execution when injected elsewhere. Existing steering
approaches, however, treat each task in isolation---vectors are extracted per
task, stored, and retrieved, leaving models unable to exploit shared structure
when a novel task arrives. We ask: what part of a new task's steering operator
can be \emph{composed} from a dictionary of previously mined skills, and how
should the composable and task-specific parts be combined? We present \method{}
(Compositional Adaptive Subspace Steering), a fully training-free framework
that (i) mines a dictionary of low-rank skill subspaces from contrastive
activations at two residual-stream depths, after projecting out a shared
task-generic component that we identify as the dominant cause of interference
in naive vector composition; (ii) given $k\le 4$ demonstrations of a novel
task, solves a joint group-sparse coding problem---with closed-form block
updates and support shared across layers---to identify the relevant skills;
and (iii) steers the model with a query-adaptive affine operator that combines
the demonstrations' denoised activation direction with a subspace correction
derived from the recovered support. We prove support-recovery guarantees under
a block-coherence condition and bound steering sub-optimality by the coding
residual, which doubles as a practical reliability signal. On 32 tasks and
Llama-3.1-8B, \method{} recovers a median 86\% of oracle task-vector
performance on held-out tasks from only 4 examples (naive composition: 0\%),
identifies the constituent skills of compound tasks with 0.68 recall,
improves monotonically with dictionary size, and amortizes to a $10\times$
per-query token saving over full in-context prompting beyond ${\sim}17$
queries per task.
\end{abstract}

\begin{CCSXML}
<ccs2012>
 <concept>
  <concept_id>10010147.10010257</concept_id>
  <concept_desc>Computing methodologies~Machine learning</concept_desc>
  <concept_significance>500</concept_significance>
 </concept>
</ccs2012>
\end{CCSXML}
\ccsdesc[500]{Computing methodologies~Machine learning}

\keywords{large language models, activation steering, task vectors, sparse
coding, representation mining}

\maketitle

\section{Introduction}
Large-scale LLM serving pipelines face a long tail of lightweight,
recurring tasks. The standard remedy---prepending $k$-shot demonstrations to
every query---multiplies prompt length by an order of magnitude and is paid
again on every call: at $10$-shot prompting, a task served a thousand times
costs ${\sim}131$k prompt tokens where the queries alone would cost
${\sim}10$k (Table~\ref{tab:cost}). A growing body of work shows that the
demonstrations' effect is substantially captured by compact structure in the
model's activation space: \emph{task vectors} or \emph{function vectors}
distilled from in-context prompts can be injected into the residual stream to
trigger task execution without any demonstrations in
context~\cite{hendel2023task,todd2024function}. Activation-level task
structure is not an artifact of one probe: different tasks induce partitioned,
clusterable hidden representations across layers~\cite{seap2025}.

Yet existing steering methods treat tasks as islands. Per-task vectors are
extracted offline and retrieved by similarity at serving
time~\cite{elicit2024}; adaptive variants train a generator network to map
demonstrations to vectors~\cite{atv2025}. Neither can serve a task absent
from the library without new training, and single-vector representations
provably degrade on tasks whose activation signature is high-rank
\cite{ranklimit2026}. The blind spot is \emph{unseen} tasks: can a steering
operator for a task never seen before be assembled, at negligible cost, from
skills mined earlier?

Naive assembly fails. Simply averaging dictionary vectors steers the model to
0\% of oracle performance in our experiments (\S\ref{sec:e1}). We trace this
failure to a structural property: every contrastively extracted task vector
shares a large \emph{task-generic} component---the signature of ``an ICL task
is being executed''---whose superposition drowns task-specific directions.
Projecting out a rank-1 shared subspace drops the median pairwise cosine
between task means from $0.34$ to $0.18$ and doubles downstream composition
accuracy (\S\ref{sec:e4}). Once removed, a second structure appears: the
residual task vectors are \emph{partially} compositional. A held-out task's
representation is well approximated inside the span of a handful of related
skills---its group-sparse code selects semantically correct neighbors
(currency$\to$\{capital, landmark, park\}; French$\to$\{German,
Spanish\})---while an irreducible task-specific component remains outside
every other task's subspace. Effective adaptation must therefore
\emph{combine} the composable part (which the dictionary can identify,
denoise, and correct toward) with the task-specific part (which only the $k$
demonstrations carry).

\method{} operationalizes this in three training-free stages
(Fig.~\ref{fig:pipeline}): dictionary mining with common-component removal
(\S\ref{sec:mining}); joint group-sparse coding across two injection layers
with closed-form block updates, millisecond solve times, and a shared support
(\S\ref{sec:coding}); and a query-adaptive affine steering operator whose
direction is the demonstrations' denoised contrastive activation and whose
correction pulls hidden states toward the recovered skill subspace
(\S\ref{sec:injection}). The coding residual $\varepsilon$ bounds steering
sub-optimality (Prop.~\ref{prop:transfer}) and doubles as a deployment-time
reliability signal: when $\varepsilon$ is large the system can fall back to
full ICL.

\textbf{Contributions.}
(1)~The first training-free framework that adapts activation steering to
unseen tasks by composing mined skill subspaces, with an explicit,
interpretable decomposition into composable and task-specific parts.
(2)~A mechanism finding: a shared task-generic activation component is the
dominant cause of interference in vector composition; removing it (at the
right rank) is necessary for compositionality to emerge, doubling composition
accuracy, with a sharply non-monotone rank ablation.
(3)~Theory: support-recovery guarantees for the joint multi-layer
group-sparse code under a block-coherence condition, and an error-transfer
bound linking coding residual to steering performance---together with
measured coherence values that delineate where the theory's conditions hold
(within-family coherence violates them exactly where we observe
wrong-neighbor steering).
(4)~A scaling and cost narrative for activation mining as a reusable asset:
composition quality grows monotonically with dictionary size, dictionaries
need only 25 contrastive samples per skill, and serving amortizes to a
$10\times$ token saving over ICL at $10^3$ queries.

\section{Related Work}
\textbf{Task and function vectors.} Hendel et al.~\cite{hendel2023task} and
Todd et al.~\cite{todd2024function} show ICL compresses task identity into
injectable vectors; follow-ups analyze their geometry and limits, including
rank limitations of single vectors~\cite{ranklimit2026}. These works extract
and analyze \emph{per-task} representations; we \emph{synthesize} operators
for unseen tasks and use subspaces as primitives, directly addressing the
high-rank failure.

\textbf{Vector libraries and adaptivity.} ELICIT~\cite{elicit2024} retrieves
from a capability library; ATV~\cite{atv2025} trains a small network to
generate task vectors. Retrieval cannot leave the library of known tasks, and
generation requires training. \method{} is training-free and returns an
explicit compositional explanation (the sparse code) with each operator.

\textbf{Activation steering.} Steering vectors add fixed directions for known
behaviors~\cite{turner2023actadd,zou2023repe}; conceptors generalize to
ellipsoidal regions and Boolean combinations of \emph{known}
behaviors~\cite{conceptor2024}. We instead \emph{infer} the combination for an
unknown task from $k$ examples; conceptor-style composition enters our
experiments as the anchor-blend ablation.

\textbf{Mining task structure in activations.} SEAP~\cite{seap2025} clusters
task activations to prune computation; contextual-sparsity methods exploit
similar structure for efficiency~\cite{dejavu2023}. We use the same raw
material---activation data across tasks---but mine it into a reusable skill
dictionary that \emph{adds} capability (steering to unseen tasks) rather than
removing computation.

\section{Skill Structure in Activation Space}
\label{sec:structure}
\textbf{Setup and notation.} Let $M$ be a decoder-only LM with $L$ layers and
residual width $d$ ($L{=}32$, $d{=}4096$ for Llama-3.1-8B). For prompt $p$,
$h_\ell(p)\in\R^d$ is the residual-stream state of the last token at layer
$\ell$. We assume a base task set $\mathcal T=\{1,\dots,T\}$ with datasets of
input--output pairs, organized in five families (knowledge mapping,
linguistic transforms, translation, algorithmic/symbolic, list selection);
$T{=}32$ in our experiments.

\textbf{Contrastive extraction.} For task $t$ and sample $i$ we build a
10-shot ICL prompt $p_i^{t,+}$ and a corrupted twin $p_i^{t,-}$ with the same
inputs but derangement-shuffled labels (format and vocabulary preserved, task
mapping destroyed), and record the differential activation
$g_i^{(t)}=h_\ell(p_i^{t,+})-h_\ell(p_i^{t,-})$ at every layer in one forward
pass. Per task we use $n{=}100$ pairs ($n{=}25$ suffices; \S\ref{sec:e4}).

\textbf{The shared task-generic component (H1).} Stacking task means
$\bar g_t$ reveals a dominant shared direction: the top left-singular vector
of $[\bar g_1 \cdots \bar g_T]$ captures format-following and
answer-emission signals common to all ICL tasks. It is the primary source of
interference when vectors are added: with it, the median pairwise
$|\cos(\bar g_s,\bar g_t)|$ is $0.34$; projecting out rank $r_0{=}1$ lowers
this to $0.18$, and composition accuracy doubles ($0.25\to0.50$,
Table~\ref{tab:ablation}). Removing too much ($r_0{\ge}8$) destroys the very
directions steering needs---the curve is sharply non-monotone, which supports
reading $U_0$ as a real, low-rank shared component rather than generic
denoising.

\section{\method}
\subsection{Module 1: Mining the Skill Dictionary}
\label{sec:mining}
Fix a small set of injection layers $\mathcal L$ ($|\mathcal L|{=}2$;
selection below). Per layer $\ell$: compute the shared component
$U_0^{(\ell)}\in\R^{d\times r_0}$ by SVD of stacked task means; project it
out of all differential activations,
$\tilde G_t^{(\ell)}=(I-U_0^{(\ell)}U_0^{(\ell)\top})G_t^{(\ell)}$; and take
the truncated SVD basis $U_t^{(\ell)}\in\R^{d\times r_{t\ell}}$ covering 90\%
spectral energy (capped at rank 16), with anchor
$\mu_t^{(\ell)}=$ mean of $\tilde G_t^{(\ell)}$. The dictionary entry for
skill $t$ stacks layers block-diagonally:
\[
\mathbf U_t=\begin{bmatrix}U_t^{(\ell_1)} & 0\\[2pt] 0 &
U_t^{(\ell_2)}\end{bmatrix}\in\R^{2d\times r_t},\qquad
\boldsymbol\mu_t=\begin{bmatrix}\mu_t^{(\ell_1)}\\[2pt]
\mu_t^{(\ell_2)}\end{bmatrix},
\]
so columns remain orthonormal and the group-sparse solver below applies
unchanged, while coefficients stay layer-specific and the \emph{support} is
shared across layers.

\textbf{Layer selection.} A one-time causal scan over single layers, followed
by a small grid over pairs using \emph{oracle} injection recovery on held-in
tasks, selects $\mathcal L=\{12,16\}$ for Llama-3.1-8B. Two depths matter
because task families peak at different layers (knowledge ${\sim}12$,
linguistic ${\sim}16{-}18$); joint injection lifts mean oracle recovery from
$0.44$ (best single layer, additive) to $0.76$ with no per-task tuning.

\subsection{Module 2: Joint Group-Sparse Coding}
\label{sec:coding}
At test time a novel task $t^\star$ provides $k\le4$ examples. For each
example $j$ we form $m{=}6$ (clean, corrupted) prompt pairs---shots are the
other $k{-}1$ examples in resampled orders, corruption re-drawn---and average
the differential activations into a denoised, stacked, deshared
representation $z_j\in\R^{2d}$; $z=\frac1k\sum_j z_j$. Averaging over prompt
resamplings is essential: it raises $\cos(z,\boldsymbol\mu_{t^\star})$ from
${\sim}0.5$ to ${\sim}0.7{-}0.8$.

We solve the group LASSO
\begin{equation}
\min_{c_1,\dots,c_T}\ \tfrac12\Big\|z-\sum_t \mathbf U_t c_t\Big\|_2^2
+\lambda\sum_t\sqrt{r_t}\,\|c_t\|_2
\label{eq:glasso}
\end{equation}
by block coordinate descent; orthonormal blocks make each update an exact
group soft-threshold, and the full solve takes milliseconds on CPU.
$\lambda$ is set without any validation data by walking the path from
$\lambda_{\max}$ downward with warm starts and keeping the best-fitting
solution whose support size is at most $s_{\max}{=}5$. Outputs: support
$S=\{t:\|c_t^*\|>0\}$, reconstruction
$\Delta_{\text{rec}}=\sum_{t\in S}\mathbf U_t c_t^*$, and residual
$\varepsilon=\|z-\Delta_{\text{rec}}\|/\|z\|$.

\subsection{Module 3: Query-Adaptive Hybrid Injection}
\label{sec:injection}
The recovered support is reliably semantic, but $\Delta_{\text{rec}}$ alone
under-steers: a held-out task's vector keeps $40{-}60\%$ of its energy
outside all other skills' subspaces, and that out-of-span component is
precisely what distinguishes, \eg, French from Spanish. \method{} therefore
injects the \emph{demonstration direction} and uses the dictionary for
\emph{structure}: with support weights
$w_t=\|c_t^*\|/\sum_{s\in S}\|c_s^*\|$, anchor
$\boldsymbol\mu_S=\sum_{t\in S}w_t\boldsymbol\mu_t$, subspace projector
$P_S$ (QR of concatenated support bases, per layer), and $z$ rescaled to the
support-anchor norm, the last-token state at each $\ell\in\mathcal L$ is
updated as
\begin{equation}
\begin{aligned}
h &\leftarrow h+\alpha(h)\big[\gamma\,\Delta^{(\ell)}
+P_S^{(\ell)}\big(\mu_S^{(\ell)}-h\big)\big],\\
\alpha(h)&=\min\!\Big(\alpha_{\max},\,
\beta\,\|(I-P_S^{(\ell)})(h-\mu_S^{(\ell)})\|/\|h\|\Big),
\end{aligned}
\label{eq:inject}
\end{equation}
with $\Delta$ the (rescaled) $z$ and $(\gamma,\beta,\alpha_{\max})=(1,2,1)$
fixed once on three held-in tasks and frozen for all experiments and models.
Intuition: the further $h$ lies from the synthesized task manifold, the harder
it is pulled; states already on-manifold are barely touched. Setting
$\Delta=\Delta_{\text{rec}}$ instead (pure synthesis) is the key ablation: it
isolates how much of a new task is dictionary-expressible
(Table~\ref{tab:ablation}).

\section{Theoretical Analysis}
\label{sec:theory}
Define block coherence
$\mu_B=\max_{s\ne t}\sigma_{\max}(\mathbf U_s^\top\mathbf U_t)$, the largest
principal-angle cosine between distinct skill subspaces.

\begin{proposition}[Support recovery]\label{prop:support}
Suppose $z=\sum_{t\in S_0}\mathbf U_t c_t^0+e$ with $|S_0|=k_0$ and
$\|e\|\le\epsilon$. If $k_0<\tfrac12(\mu_B^{-1}+1)$ and $\lambda$ is chosen
appropriately for $\epsilon$, the group-LASSO solution satisfies
$\hat S\subseteq S_0$ with coefficient error $O(\epsilon)$.
\end{proposition}
The proof instantiates block-sparse recovery results for unions of
subspaces~\cite{eldar2009robust,eldar2010block} and group-LASSO
consistency~\cite{yuan2006model}; we contribute the \emph{verification}:
common-component removal is what brings measured coherence into a workable
regime (median pairwise $\mu_B$ vs.\ $r_0$; Table~\ref{tab:mub}), and the
condition's failure is diagnostic---the largest coherences occur inside the
translation family ($\mu_B$ up to $0.95$), exactly where we observe correct
\emph{family} recovery but wrong-member steering (\S\ref{sec:analysis}).

\begin{proposition}[Error transfer]\label{prop:transfer}
If task performance is locally linear in the injected state,
$m(h)=w^\top h+O(\|h\|_{loc}^2)$, then
$|m(h+\Delta)-m(h+\Delta^{or})|\le\|w\|\,(\varepsilon\|z\|+\delta_{est})$,
where $\delta_{est}$ is the finite-sample subspace estimation error controlled
by Davis--Kahan at rate $O(\sigma/(\sqrt n\,\mathrm{gap}))$.
\end{proposition}
Empirically $\varepsilon$ anti-correlates with recovery (Fig.~\ref{fig:eps}),
grounding its use as a deployment-time fallback trigger.

\section{Experiments}
\label{sec:exp}
\textbf{Setup.} 32 tasks from the function-vectors
benchmark~\cite{todd2024function} across five families; per task, 50 held-out
evaluation queries disjoint from dictionary and demonstration pools; 5 seeds
for demonstration sampling. Main model Llama-3.1-8B-Instruct; cross-model
validation on Qwen3-4B and Llama-3.2-3B. Metric: strict first-words exact
match (case-sensitive for compound tasks). Recovery
$\rho=(acc_{\text{syn}}-acc_{\text{zs}})/(acc_{\text{or}}-acc_{\text{zs}})$
against each task's own-subspace oracle injection; tasks with oracle-zs gap
$<5$pts are excluded from $\rho$ aggregation and reported separately. A
single RTX~3090 suffices for the entire study (${\sim}40$ GPU-hours).

\subsection{E1: Unseen-Task Synthesis (Leave-One-Task-Out)}
\label{sec:e1}
For each held-out task the dictionary (including $U_0$) is rebuilt from the
remaining 31 tasks. With $k{=}4$ examples, \method{} attains median
$\rho=\mathbf{0.86}$; 70\% of the 27 valid tasks exceed $\rho{=}0.6$
(Fig.~\ref{fig:e1}). Median $\rho$ is $0.37$ at $k{=}1$ and already $0.89$ at
$k{=}2$. By family (k=4): linguistic $1.01$, selection $0.93$, knowledge
$0.84$, algorithmic $0.65$, translation $0.43$. Baselines: naive anchor
averaging $\rho=0.00$; raw demonstration vector without the dictionary
(``$z$-only'') trails \method{} by $3.9$pts mean accuracy at $k{=}4$; pure
reconstruction trails by $16.8$pts, quantifying the irreducible task-specific
component.

\subsection{E2: Compound Tasks and Interpretability}
Ten compound tasks chain two dictionary skills (\eg\
\texttt{singular-plural}$\circ$\texttt{capitalize}); none is in the
dictionary, and scoring is case-sensitive so the second component is only
credited when actually executed. The sparse code recovers constituent skills
with recall $0.68$ (precision $0.36$; coefficient heatmap in
Fig.~\ref{fig:heatmap}). Steering executes the composition where the
components' formats are compatible (\texttt{singular-plural+capitalize}
$0.86$, \texttt{present-past+capitalize} $0.50$) and \method{} is the only
non-trivial method overall ($0.16$ vs.\ retrieval $0.03$, naive component
averaging $0.05$; 10-shot ICL $0.61$): retrieval cannot compose, and its
failure here is categorical, not marginal.

\subsection{E4: Ablations}
\label{sec:e4}
Table~\ref{tab:ablation} summarizes (8 representative tasks, $k{=}4$, 3
seeds). Common-component removal is load-bearing: $r_0{=}0/1/2/8/16 \to$
mean accuracy $0.25/0.50/0.42/0.17/0.02$. Injection layers:
$\{12{,}16\}$ ($0.42$) $>$ $\{12{,}14{,}16\}$ ($0.40$) $>$ single 16 ($0.28$)
$>$ single 12 ($0.14$). Direction source: hybrid $0.42$ $>$ $z$-only $0.41$
$\gg$ reconstruction $0.10$ $>$ anchor blend $0.05$. Solver: group LASSO
$\approx$ block-OMP $>$ simplex $>$ least squares. Dictionary build cost:
$n{=}25$ samples per task matches $n{=}100$.

\subsection{E5: Dictionary Scaling and Serving Cost}
Accuracy on six fixed held-out tasks grows monotonically with dictionary
size---$0.14/0.44/0.51/0.52$ at $T'{=}5/10/15/31$ (mean over 5 random
sub-dictionaries)---supporting the mined-dictionary-as-asset framing: every
added skill improves zero-training transfer to unseen tasks
(Fig.~\ref{fig:scale}). Serving cost (Table~\ref{tab:cost}): a 10-shot ICL
prompt averages $131$ tokens vs.\ $10.5$ zero-shot; \method's one-time
extraction costs $2{,}211$ tokens, breaking even at ${\sim}17$ queries and
saving $10.2\times$ at $10^3$ queries per task.

\subsection{E6: Cross-Model Validation}
[E6 running: Qwen3-4B and Llama-3.2-3B with depth-proportional layer pairs;
numbers to be inserted.]

\subsection{Failure Analysis}
\label{sec:analysis}
Stored generations let us separate failure modes. (i)~\emph{Format
non-compliance}: on \texttt{person-occupation}/\texttt{person-sport} the
steered model emits biographies containing the answer (contains-metric
$0.29$ vs.\ exact $0.02$; oracle contains $0.38$)---steering conveys the
mapping but not the answer-only format. (ii)~\emph{Coherent-neighbor
attraction}: held-out \texttt{english-french} recovers the translation
family but is pulled toward Spanish/German anchors by the correction term,
producing fluent wrong-language answers; this is the $\mu_B$ condition of
Prop.~\ref{prop:support} failing within-family, observable in advance from
the dictionary itself. (iii)~$\varepsilon$ predicts $\rho$
(Fig.~\ref{fig:eps}), enabling selective fallback to full ICL.

\section{Conclusion}
\method{} shows that the activation footprint of an unseen task decomposes
into a part that is compositional over previously mined skills and an
irreducible task-specific remainder---and that a training-free system
exploiting both recovers most of oracle steering performance from four
examples at a fraction of ICL's serving cost. Mining activation structure
once and reusing it across the task long tail turns steering from a
per-task artifact into a growing knowledge asset.

\bibliographystyle{ACM-Reference-Format}
\bibliography{references}

\appendix
\section{Reproducibility}
Code, task splits, frozen hyperparameters, and per-run generations are
released; all experiments run on one 24GB GPU. [Proofs, full per-task
tables, and additional ablations to be inserted.]

\end{document}
